% Generated draft: requires mathematical and source review.
\documentclass[11pt]{article}
\usepackage[margin=25mm]{geometry}
\usepackage{fontspec}
\setmainfont{Latin Modern Roman}
\usepackage{amsmath,amssymb,mathtools,unicode-math}
\setmathfont{Latin Modern Math}
\usepackage{enumitem,needspace,xurl,hyperref}
\setlength{\parindent}{0pt}
\setlength{\parskip}{5pt}
\setlength{\emergencystretch}{3em}
\newcommand{\N}{\mathbb{N}}
\newcommand{\Z}{\mathbb{Z}}
\newcommand{\Q}{\mathbb{Q}}
\newcommand{\R}{\mathbb{R}}
\newcommand{\C}{\mathbb{C}}
\newcommand{\F}{\mathbb{F}}
\newcommand{\A}{\mathbb{A}}
\newcommand{\PP}{\mathbb{P}}
\newcommand{\E}{\mathbb{E}}
\newcommand{\eps}{\varepsilon}
\newcommand{\dd}{\,\mathrm{d}}
\begin{document}
\section*{UnsolvedMath ID: 800116}
\textbf{Source problem code:} catalogue entry
\subsection{Definitions and mathematical statement}
Let \(M\) be a real analytic manifold equipped with a real analytic foliation \(\mathcal F\) of codimension \(1\). A subset is saturated if it is a union of leaves. A minimal set is a nonempty closed saturated subset containing no proper nonempty closed saturated subset. Conventionally, an exceptional minimal set is a nowhere dense minimal set that is not a single leaf; the source does not specify its ambient hypotheses or its precise convention for this term.

Let \(X\) be such an exceptional minimal set. A local transversal \(T\) is an interval transverse to the leaves. An extremal leaf of \(X\), in the sense described in the supplied record, is a leaf through a point of \(X\cap T\) that is an endpoint of a component of \(T\setminus(X\cap T)\). Here the endpoint must lie inside the transversal, so that an artificial boundary of the chosen transversal does not create an extremal leaf. Write \(\mathcal E(\mathcal F,X)\) for the collection of these leaves.

The ends of a leaf \(L\) are taken in its intrinsic manifold topology. For a compact exhaustion \(K_1\subseteq K_2\subseteq\cdots\) of \(L\), their space is
\[
\operatorname{Ends}(L)=\varprojlim_j\pi_0(L\setminus K_j),
\]
where \(\pi_0\) denotes connected components and the bonding maps are induced by inclusion.

The supplied record states that, for real analytic codimension-one foliations, the presence of an exceptional minimal set implies that an extremal leaf has infinitely many ends. The wording does not explicitly distinguish the following quantifiers:
\[
\forall L\in\mathcal E(\mathcal F,X),\qquad
\bigl|\operatorname{Ends}(L)\bigr|\geq\aleph_0,
\]
or
\[
\exists L\in\mathcal E(\mathcal F,X)\quad
\bigl|\operatorname{Ends}(L)\bigr|\geq\aleph_0.
\]
These are alternative readings of the supplied assertion, not two independently established claims. Compactness, boundary, connectedness, orientability, transverse orientability, and any dimension restrictions on \(M\) are not specified in the record.

The source question asks whether the preceding theory for analytic group actions can be generalized to real analytic codimension-one foliations, and whether such a generalization can establish the Dippolito conjecture for foliations. The referenced theory and the statement of that conjecture are absent from the supplied record. Consequently, a precise extension theorem or conjectural conclusion cannot be specified from this evidence alone.
\subsection{Short English statement}
Question 57 asks whether a theory of analytic group actions extends to real analytic codimension-one foliations and can establish the Dippolito conjecture. It supplies a claim relating exceptional minimal sets to extremal leaves with infinitely many ends, but leaves the quantifier, ambient hypotheses, referenced theory, and conjecture statement unresolved.
\subsection{Sources}
\begin{itemize}
\item Supplied local UnsolvedMath registry record: permanent ID \(800116\), exact source problem code \texttt{catalogue entry}, title ``Foliations and Diffeomorphism Groups --- Question 57.'' The registry attributes the problem to Nariman and Yazdi, with proposed year \(2025\). Its status is supplied as \texttt{open}; this label is not an independently verified mathematical assessment.
\item Nariman and Yazdi, \emph{Problem List on Foliations and Diffeomorphism Groups} (2025), Question 57. Supplied source URL: \url{https://arxiv.org/abs/2503.18274}. The record reports access on 2026-10-16 and extraction by \texttt{arxiv\_latex}; these are supplied provenance details, not evidence of independent source consultation for this draft.
\item Supplied upstream revision: \texttt{d371ee917ac9647fbb06cd6fc529deec6f3dda82}. The record gives no separate reference for the Dippolito conjecture or for the analytic group-action theory to which the question refers.
\end{itemize}
\end{document}
